\documentclass[11pt,a4paper]{article}
\usepackage{gtu-el-yazisi}
\begin{document}
\notbilgi{Alan Teorisi}{Uygulama 4 -- Yük sistemlerinde potansiyel ve enerji}{Uygulama4.pdf}

\sayfa{1}
\noindent\underline{Sayfa 27.Pr11}

\elcizimi[0.8]{sekiller/uygulama-notlari-Uygulama4-s1-1.png}[1]

\formulkutu{V=\frac{Q}{4\pi\varepsilon_0\,r}}

a)
\begin{align*}
  V_C &= \frac{Q}{4\pi\varepsilon_0(x_0+a)}+\frac{Q}{4\pi\varepsilon_0(a-x_0)} \\
      &= \frac{Q}{4\pi\varepsilon_0}\Bigl[\frac{1}{a+x_0}+\frac{1}{a-x_0}\Bigr] \\
      &= \frac{Q}{4\pi\varepsilon_0}\,\frac{2a}{a^2-x_0^2} \\
  V_C &= \frac{Q}{2\pi\varepsilon_0}\,\frac{a}{a^2-x_0^2}
\end{align*}

\formulkutu{\text{potansiyel enerji}\quad W_C=qV_C=\frac{qQa}{2\pi\varepsilon_0}\,\frac{1}{a^2-x_0^2}}

b) $q$ yükü serbest bırakılırsa ne olur
\[ W_T=\overbrace{\underbrace{W_{P_C}}_{\text{C noktasında potansiyel enerji}}}^{\text{potansiyel enerji}}+\overbrace{W_{K_C}}^{0}\ \to\ \text{cisim hareketsiz} \]
\[ W_T=\frac{qQa}{2\pi\varepsilon_0}\,\frac{1}{a^2-x_0^2} \]
\[ W_T=W_{P_x}+W_{K_c}=\frac{qQa}{2\pi\varepsilon_0}\,\frac{1}{a^2-x^2}+\frac{1}{2}mv^2 \]
\[ =\frac{qQa}{2\pi\varepsilon_0}\,\frac{1}{a^2-x_0^2} \]
\[ v=\sqrt{\frac{aqQ}{m\pi\varepsilon_0}\Bigl(\frac{1}{a^2-x_0^2}-\frac{1}{a^2-x^2}\Bigr)} \]

\sayfa{2}
\noindent c) Bu hızın 0 olduğu nokta

\medskip
\noindent $x=x_0$ \ ve $x=-x_0$ \quad $v=0$ olur \quad yük $-x_0$, $x_0$ arasında harmonic hareket yapar

\medskip
\noindent $x=0$ iken hız maksimumdur

\bigskip
\noindent d) başlangıçta serbest bırakılmaz ve $v_0$ hızıyla bırakılırsa sonuç nasıl değişir

\medskip
\noindent C noktasında
\[ W_T=W_{P_C}+W_{K_C} \]
\formulkutu{=\frac{qQa}{2\pi\varepsilon_0}\,\frac{1}{a^2-x_0^2}+\frac{1}{2}mv_0^2}

\noindent Herhangi bir $x$ noktasında
\[ W_T=\frac{qQa}{2\pi\varepsilon_0}\,\frac{1}{a^2-x^2}+\frac{1}{2}mv^2= \]
\[ =\frac{qQa}{2\pi\varepsilon_0}\,\frac{1}{a^2-x_0^2}+\frac{1}{2}mv_0^2 \]
(Kırmızı kalemle iki satırda da $2\pi\varepsilon_0$ ve $\tfrac{1}{2}$ içindeki $2$'ler çizilmiş.)
\[ \Rightarrow\ v^2=\frac{aQq}{m\pi\varepsilon_0}\,\frac{1}{a^2-x_0^2}-\frac{qQa}{m\pi\varepsilon_0}\,\frac{1}{a^2-x^2}+v_0^2 \]
\[ v=\sqrt{\frac{qQa}{m\pi\varepsilon_0}\Bigl(\frac{1}{a^2-x_0^2}-\frac{1}{a^2-x^2}\Bigr)+v_0^2} \]

\noindent salındığı noktalar $x_1$ ve $x_{-1}$ olsun
\begin{align*}
  v(x_1) &= 0 \\
  v(-x_1) &= 0
\end{align*}
\[ \frac{qQa}{m\pi\varepsilon_0}\Bigl(\frac{1}{a^2-x_0^2}-\frac{1}{a^2-x_1^2}\Bigr)+v_0^2=0 \]

\sayfa{3}
\[ \Rightarrow\ \frac{1}{a^2-x_1^2}=\frac{1}{a^2-x_0^2}+\frac{m\pi\varepsilon_0}{aqQ}\,v_0^2 \]
\[ x_1=\pm\sqrt{a^2-\frac{1}{\dfrac{1}{a^2-x_0^2}+\dfrac{m\pi\varepsilon_0 v_0^2}{aqQ}}} \]

\noindent $x_1$ her zaman $x_0$'dan her zaman büyüktür.

\medskip
\noindent{\color{kirmizikalem}\underline{S.27 Pr12}}

\elcizimi[0.4]{sekiller/uygulama-notlari-Uygulama4-s3-1.png}[3]

\noindent a) $P(x,0,0)$ noktasında \quad $V=?$ \ $\vec E=??$
\[ V(x,0,0)=\frac{Q}{4\pi\varepsilon_0}\,\frac{1}{\sqrt{x^2+h^2}}+\frac{Q_1}{4\pi\varepsilon_0}\,\frac{1}{\sqrt{x^2+h^2}} \]
\[ \vec E=-\operatorname{grad}V=-\frac{\partial}{\partial x}\Bigl(\frac{2Q}{4\pi\varepsilon_0}\,\frac{1}{\sqrt{x^2+h^2}}\Bigr)\vec e_x \]
\formulkutu{\vec E=\frac{Q}{2\pi\varepsilon_0}\,\frac{x}{(x^2+h^2)^{3/2}}\,\vec e_x}

\noindent b) $P(x_0,0,0)$ noktasına $-q$ yükü yerleştiriliyor. Bu noktanın sahip olduğu potansiyel enerji ve üzerine etki eden kuvvet nedir?
\begin{align*}
  W_{P_0} &= -q\,V(x_0,0,0) = -\frac{qQ}{2\pi\varepsilon_0}\,\frac{1}{\sqrt{x_0^2+h^2}} \\
  \vec F &= -q\,\vec E(x_0,0,0) = -\frac{qQ}{2\pi\varepsilon_0}\,\frac{x_0}{(x_0^2+h^2)^{3/2}}\,\vec e_x
\end{align*}

\sayfa{4}
\noindent c) $-q$ yükü $v_0=0$ ile serbest bırakılıyor. Bu yükün $(-a,0,0)$ ve $(a,0,0)$ arasında salınım yaptığını gösteriniz
\[ W_T=W_{P_{x_0}}+\overbrace{W_{K_{x_0}}}^{0}=\frac{-qQ}{2\pi\varepsilon_0\sqrt{x_0^2+h^2}} \]

\noindent herhangi bir $x$ noktasında
\[ W_T=W_{P_x}+\frac{1}{2}mv^2 \]
\[ \frac{-qQ}{2\pi\varepsilon_0\sqrt{x_0^2+h^2}}=\frac{-qQ}{2\pi\varepsilon_0\sqrt{x^2+h^2}}+\frac{1}{2}mv^2 \]
(Kırmızı kalemle $2\pi\varepsilon_0$ ve $\tfrac{1}{2}$ içindeki $2$'ler çizilmiş.) \qquad $a=x_0$
\[ v=\pm\sqrt{\frac{qQ}{m\pi\varepsilon_0}\Bigl(\frac{1}{\sqrt{x^2+h^2}}-\frac{1}{\sqrt{x_0^2+h^2}}\Bigr)} \]
\begin{align*}
  x &= x_0 \\
  x &= -x_0
\end{align*}

\noindent d) $\vec v_0=v_0\,\vec e_x$ olsun

\noindent $v_0$ pozitif veya negatiftir

\noindent yükün $(-x_1,0,0)$ ve $(x_1,0,0)$ arasında salınım yaptığını gösteriniz, $x_1$'i bulunuz

\noindent $x_1>a$ olduğunu gösteriniz.
\[ W_T=W_{P_{x_0}}+\frac{1}{2}mv_0^2=\frac{-qQ}{2\pi\varepsilon_0\sqrt{x_0^2+h^2}}+\frac{1}{2}mv_0^2 \]
\noindent herhangi bir $x$ noktasında
\[ W_T=\frac{-qQ}{2\pi\varepsilon_0\sqrt{x^2+h^2}}+\frac{1}{2}mv^2 \]

\sayfa{5}
\[ =\frac{-qQ}{2\pi\varepsilon_0\sqrt{x_0^2+h^2}}+\frac{1}{2}mv_0^2 \quad\text{($x_0$'daki)} \]
\[ \Rightarrow\ v=\sqrt{\frac{qQ}{m\pi\varepsilon_0}\Bigl(\frac{1}{\sqrt{x^2+h^2}}-\frac{1}{\sqrt{x_0^2+h^2}}\Bigr)+v_0^2} \]

\noindent $v(-x_1)=0$ ve $v(x_1)=0$ olduğunu gösteriniz.
\[ \sqrt{\frac{qQ}{m\pi\varepsilon_0}\Bigl(\frac{1}{\sqrt{x_1^2+h^2}}-\frac{1}{\sqrt{x_0^2+h^2}}\Bigr)+v_0^2}=0 \]
\[ \frac{qQ}{m\pi\varepsilon_0}\,\frac{1}{\sqrt{x_1^2+h^2}}-\frac{qQ}{m\pi\varepsilon_0}\,\frac{1}{\sqrt{x_0^2+h^2}}+v_0^2=0 \]
\[ \Rightarrow\ \frac{1}{\sqrt{x_1^2+h^2}}=\frac{1}{\sqrt{x_0^2+h^2}}-\frac{m\pi\varepsilon_0}{qQ}\,v_0^2 \]
\[ x_1^2+h^2=\frac{1}{\Bigl(\dfrac{1}{\sqrt{x_0^2+h^2}}-\dfrac{m\pi\varepsilon_0 v_0^2}{qQ}\Bigr)^2} \]
\[ \Rightarrow\ x_1=\pm\sqrt{\frac{1}{\Bigl(\dfrac{1}{\sqrt{x_0^2+h^2}}-\dfrac{m\pi\varepsilon_0 v_0^2}{qQ}\Bigr)^2}-h^2} \]
\[ x_1>a=x_0 \]

\sayfa{6}
\elcizimi[0.35]{sekiller/uygulama-notlari-Uygulama4-s6-1.png}[6]

\noindent a) $V=?$

\noindent b) $\vec E$

\noindent c) $V(r)$ ve $\vec E(r)$ r'ye bağımlı değişimi nedir
\[ Q=\int_{\theta=0}^{\pi}\int_{\phi=0}^{2\pi}\rho_s\,a^2\sin\theta\,d\phi\,d\theta \]
\[ Q=4\pi a^2\rho_s\ \Rightarrow\ \rho_s=\frac{Q}{4\pi a^2} \]
\[ dV=\frac{\rho_s\,ds}{4\pi\varepsilon_0\,r'} \]
\[ V=\int_{\theta=0}^{\pi}\int_{\phi=0}^{2\pi}dV=\int_{\theta=0}^{\pi}\int_{0}^{2\pi}\frac{\rho_s\,ds}{4\pi\varepsilon_0\,r'} \]
\begin{align*}
  ds &= a^2\sin\theta\,d\phi\,d\theta \\
  \rho_s &= \frac{Q}{4\pi a^2}
\end{align*}
\[ \vec E=-\operatorname{grad}V \]

\end{document}
